\section{Half-Euler extraction and polynomial bit complexity}\label{sec:extraction}
\begin{lemma}[Half-Euler walk]\label{lem:half-euler}
Let a finite nonempty balanced directed multigraph have a reversing
involution on vertices and edge copies, no fixed edge, connected used
support, and a fixed root in that used support. If every fixed vertex has even outdegree,
there is a directed closed walk from the root using exactly one edge
from each involution pair. It can be constructed in polynomial time.
\end{lemma}
\begin{proof}
Quotient vertices and edges by the involution, retaining both incidences
of each edge, including loops. At a nonfixed vertex orbit choose a
representative $v$. Label an incidence positive if its compatible lift
is outgoing at $v$, and negative if incoming there. The numbers of
positive and negative incidences equal $\outdeg(v)$ and $\indeg(v)$,
so they agree. At a fixed vertex the quotient degree equals its directed
outdegree and is even.

Pair opposite signs at each nonfixed orbit and pair arbitrary incidences
at fixed vertices. The resulting transition system decomposes the
quotient edges into closed alternating trails. These lift consistently:
at a nonfixed orbit an arrival through a positive incidence puts the
lift at the mate of $v$, where departure requires a negative incidence;
the other case is reversed. At a fixed vertex either state is available.
The final/initial paired incidences give closure. If several trails
remain, connectedness supplies two sharing a vertex. Cross-switch two
incidence pairs there, aligning opposite signs at a nonfixed orbit,
to merge the trails. This preserves all local conditions. Repeating
produces a quotient Euler trail, which lifts from the fixed root and
closes there, selecting one copy from each orbit.

This is the signed-incidence Euler construction underlying the
antecedents in Section~\ref{sec:antecedents}, with fixed vertices handled
explicitly. The general lemma allows quotient loops. Our hierarchical
graph has none, because an edge changes word length and $\rho$
preserves it.
\end{proof}

\begin{proof}[Proof of Theorem~\ref{thm:main}]
The empty cases were handled before the construction. Otherwise apply
Lemma~\ref{lem:half-euler} to $D$ from Theorem~\ref{thm:connection}, rooted at
$\eps$. Since $D$ has $U$ up and $U$ down edges and $\tau$ exchanges
them, it has $U$ edge orbits. The selected closed walk has $U$ edges;
half are up, since word length returns to zero. It visits a representative
of every used vertex orbit, and thus one orientation of every required
word. Append the letter of each up step to the output. At every moment
the current vertex is a suffix of the output so far: up appends a letter
to both, and down deletes the vertex's first letter. This is the suffix
invariant of \src{Section 7}. Therefore the word is feasible and
\begin{equation}\label{eq:final}
|T|=U/2\le W\le2\OPT(\mathcal S).
\end{equation}

For completeness we detail the finite realization and bit costs.
All finite choices use encoded-symbol order and stable IDs, including
orbits of equal least period, tied canonical maxima, transitions,
records, host anchors, threshold choices, and trail splices.
Unused alphabet symbols may be ignored after reading and checking the
explicit alphabet and involution; keep every used symbol and its mate.

There are $O(L^2)$ substring vertices of length at most $L$, and all
counts are at most $2L$. Enumerating $v,p,k$ as in Section~\ref{sec:blocking}
gives a polynomial list (the coarse bound $O(L^4)$ suffices).
Evaluate blocking by enumerating every bracket word and position.
Self-reversal is tested by comparing a primitive period word with
rotations of its reverse complement. These are polynomial operations
on explicit strings and polynomial-bit integers.

The base graph has $2W\le4L$ edge copies. Paired transitions and walk
lifts are finite scans. Least periods, normalized passages, and their
exit order statistics use explicit comparisons and sorting. The total
number of primitive phases is $W$, so storing one period of each
$z_i$ and its text is polynomial. Both axis crossings and the host
splices are scans through these stored paths.

Equality or order comparison of two periodic texts needs at most the
sum of their periods, by the agreement bound. Normalize record starts
to a turn, translating child coordinates together with the record.
All actual windows have length at most $L$, and the resulting integers
have polynomial bit length. There are $2W$ actual window visits per
turn across all layers. To find a different-alignment record for $w$,
enumerate them and each placement of $w$ within them, at most $2WL$
candidates. Align the placement to its prescribed occurrence interval;
reject identical full texts or windows not matching the source throughout.
The count and matching-window arguments prove that a candidate remains.

Threshold searches, the individual tests, batch reclassification, host
record order, and functional-graph cycle detection all operate on
polynomially many stored objects. Collective intersections scan at most
$K\le W$ start coordinates. Ordered-window paths are materialized by
prefixes, suffixes, and single letters. The paired budget proof bounds
the final graph by $2U\le4W\le8L$ edge copies, including all additions.
The quotient traversal and output take polynomial time in this explicit
graph. Comparisons and output of variable-width symbol labels are
charged their encoded bit lengths; no unbounded-label unit-cost operation,
random choice, or infinite-text oracle is used. This establishes the
claimed deterministic polynomial bit complexity at the full scope.
\end{proof}
